\section{Empirical study: tool access and decision outcomes}
\label{sec:autonomous}
We evaluate how autonomous tool access affects portfolio decisions in open-weight
language models under sequential multi-asset tasks.

\subsection{Experimental setup}
We test Qwen2.5-Coder 7B and 14B under four conditions: historical prices alone; prices
with static skill excerpts; prices with precomputed tool outputs; and autonomous access
to the library's algorithm catalog. Each condition runs across inference seeds 11 and 23,
with a budget of five interaction turns per decision (capped at 256 generated tokens per
response). Conditions run in randomized order.

In the autonomous condition, the agent receives descriptions of available tools. It can
inspect algorithm specifications, run an algorithm with custom arguments, or query a
rule-based strategy adviser. Tool calls are dispatched against trailing market snapshots
and return structured numerical results. In our evaluation, agents access 18
price-compatible quantitative routines.

A decision without valid portfolio weights leaves existing allocations unchanged.
Requests, outputs, and intermediate reasoning traces are logged separately, and final
weights are frozen before scoring.

\subsection{Data and evaluation protocol}
We evaluate on daily exchange rates for EUR against USD, JPY, GBP, and CHF from ECB
reference data. Messages use anonymized asset identifiers and relative indices rather
than calendar dates.

Each decision observes 126 trailing days. The evaluation spans the first 84 trading days
of 2025 (ending May 2), with four rebalancing points spaced 21 days apart. Decisions at
step $t$ execute at day $t+1$ prices to prevent look-ahead bias. Portfolios are long-only
without leverage, with unallocated capital held in cash. We charge 5 basis points per
side for rebalancing and final liquidation. Cash, equal-weight buy-and-hold, and
deterministic rules serve as reference baselines.

\subsection{Observed execution and returns}
Across 16 trials, models generated 64 final portfolio allocations. Autonomous 7B made 13
successful tool calls, while autonomous 14B made 30. Both models recovered from minor
JSON formatting errors within their turn budget, completing all decisions without
runtime crashes.

Tool choices diverged across model scales: 7B selected cross-sectional momentum,
Bollinger reversion, and volatility-targeted momentum, whereas 14B predominantly used
Bollinger reversion alongside trend-following indicators.

\begin{table}[t]
\centering\small
\begin{tabular}{lrr}
\toprule
Condition & 7B & 14B \\
\midrule
No library & $+0.100$ & $-3.324$ \\
Skill text & $-2.930$ & $-4.567$ \\
Fixed tool outputs & $-3.075$ & $-3.030$ \\
Autonomous library & $-3.984$ & $-3.783$ \\
\bottomrule
\end{tabular}
\caption{Mean cumulative proxy return (\%) across inference seeds with 5 bps transaction costs per side.}
\label{tab:autonomous}
\end{table}

Table~\ref{tab:autonomous} shows returns across all arms. While autonomous 14B improved
over its text-only condition by 0.783 percentage points, neither autonomous model beat
the unassisted baseline (+0.100\% for 7B, $-3.324\%$ for 14B).

Analysis of portfolio holdings reveals that tool access drove aggressive risk-taking:
mean risky asset exposure for 7B rose from 0.500 without the library to 0.969 with
autonomous tools; for 14B, exposure shifted from 1.000 (all-in) to 0.500. Autonomous
runs used 38 model calls per model across eight decisions, totaling 1,155,782 prompt
tokens and 10,165 completion tokens.

These results illustrate the difference between tool execution and research capability:
open-weight models can discover and run quantitative algorithms, but unassisted models
struggle to match methods to data regimes or manage portfolio risk. This motivates the
context pruning and memory mechanisms evaluated in Section~\ref{sec:capability_eval}.
